\section*{Chapitre \ref{ch:dv:parametrising-the-surface} --- Paramétrer la surface}

\begin{solution}{exo:dv:parametrising-the-surface:1}
$w(0)=a+b\bigl(-\rho m+\sqrt{m^2+\varsigma^2}\bigr)=0.00670$~; la volatilité implicite vaut
$\sqrt{0.00670/(91/365)}=16.4\%$.
\end{solution}

\begin{solution}{exo:dv:parametrising-the-surface:2}
À gauche $b(1-\rho)=0.062$, à droite $b(1+\rho)=0.014$. À partir de $\tilde q=1/(2\beta_L)+\beta_L/8-1/2$ avec
$\beta_L=0.062$~: $\tilde q=7.6$, donc les moments négatifs $\E S_T^{-q}$ sont finis jusqu'à environ $q=7.6$~; bien
loin de la borne 2, ce qui reviendrait à dire qu'aucun moment négatif n'est fini.
\end{solution}

\begin{solution}{exo:dv:parametrising-the-surface:3}
$0.5^2\times12/365-0.3^2\times12/365=0.00526$~; écart-type 7.3~\%~; déplacement absolu attendu
$7.3\%\times\sqrt{2/\pi}=5.8\%$.
\end{solution}

\begin{solution}{exo:dv:parametrising-the-surface:4}
L'événement ajoute de la variance à une date, et non à un taux. L'interpolation linéaire la répartit sur l'intervalle,
si bien qu'une option dont l'échéance tombe la veille de la publication se voit imputer une partie de la \omterm{def:dv:parametrising-the-surface:event}{variance d'événement} qu'elle ne subira
jamais~: avec les chiffres du chapitre, une option à 7 jours afficherait une volatilité de 49.5~\% au lieu de 32~\%.
\end{solution}

\begin{solution}{exo:dv:parametrising-the-surface:5}
$\varphi=1.095\times0.00183^{-0.478}$~; $\theta\varphi(1+|\rho|)=0.064<4$ et
$\theta\varphi^2(1+|\rho|)=1.42\le4$~: les deux sont vérifiées.
\end{solution}

\begin{solution}{exo:dv:parametrising-the-surface:6}
$\sigma_c=0.25\sqrt{(5/252)/(7/365)}=25.43\%$.
\end{solution}

\begin{solution}{exo:dv:parametrising-the-surface:7}
Sur ces cotations, les deux ajustements se situent à l'intérieur des 27 fourchettes, et leurs ailes ne diffèrent pas de plus de $2\times10^{-4}$ l'une de
l'autre (0.0623 et 0.0140 contre 0.0621 et 0.0142). L'ajustement aux fourchettes compte lorsque les cotations sont larges ou
déséquilibrées, et dans les ailes, où les cours médians sont les moins fiables.
\end{solution}

\begin{solution}{exo:dv:parametrising-the-surface:8}
Les cours médians ne sont pas les fair values~: ils portent un bruit égal à une fraction de la fourchette. Faire passer la courbe par eux
reproduit ce bruit (la spline s'écarte au plus de 0.17 point de volatilité du vrai smile, SVI de 0.08), l'injecte dans
la densité et les \omterm{def:dv:greeks-and-the-hedging-pnl:greeks}{grecques}, et ne dit rien des ailes. La bonne mesure d'erreur est la distance à la
fourchette acheteur--vendeur, que les deux respectent, et la stabilité de la surface d'un recalibrage au suivant.
\end{solution}

\begin{solution}{pb:dv:parametrising-the-surface:1}
\textbf{1.} $0.32^2\times4/365=0.001122$~; $0.561^2\times11/365=0.009485$.
\textbf{2.} $0.32^2\times11/365=0.003086$.
\textbf{3.} $0.009485-0.003086=0.0064$.
\textbf{4.} 8.0~\% et 6.4~\%.
\textbf{5.} L'échéance à 4 jours est la seule à se terminer avant la publication~; les échéances ultérieures contiennent toutes
l'événement.
\textbf{6.} 32.0~\%.
\textbf{7.} 32.1~\%.
\textbf{8.} Un second événement survenant avant ces échéances, ou un marché qui anticipe une volatilité diffusive plus élevée
après la publication.
\textbf{9.} Environ 32~\%~: la \omterm{def:dv:parametrising-the-surface:event}{variance d'événement} sort de la variance totale de l'échéance une fois la publication passée.
\textbf{10.} Le straddle a été évalué avec la \omterm{def:dv:parametrising-the-surface:event}{variance d'événement}~; un mouvement de 6~\% est inférieur à l'écart-type de 8~\%
implicite, et l'effondrement de la volatilité après la publication emporte avec lui la valeur temps.
\textbf{11.} Linéairement en variance totale, sur une horloge qui place la \omterm{def:dv:parametrising-the-surface:event}{variance d'événement} au jour 7 et la
variance diffusive sur les autres jours ouvrés.
\textbf{12.} Les échéances juste après la publication~: un saut rend le smile plus courbé aux échéances courtes,
et d'autant plus que la \omterm{def:dv:parametrising-the-surface:event}{variance d'événement} est grande par rapport à la variance diffusive.
\textbf{13.} Il faut attribuer aux jours de week-end peu ou pas de variance en \omterm{def:dv:parametrising-the-surface:business}{temps ouvré}.
\textbf{14.} La volatilité diffusive, soit 32~\%.
\textbf{15.} Le niveau de la tranche baisse fortement et la courbure ajoutée par le saut au voisinage de la monnaie disparaît~;
les paramètres du recalibrage sautent en conséquence, ce qui explique que les variations de paramètres aux dates d'événement soient attendues, et non
des alertes.
\textbf{16.} Non~: il s'agit du mouvement évalué sous la mesure de valorisation, y compris une prime pour porter le
risque d'événement.
\textbf{17.} Comparer le déplacement absolu attendu implicite à la moyenne des déplacements absolus de l'action lors de ses
publications passées, mesurés de la même façon (clôture avant contre clôture après).
\textbf{18.} Vendre le straddle à 11 jours et acheter celui à 4 jours ou une échéance ultérieure pour couvrir la volatilité
diffusive~; le risque est un mouvement supérieur à celui qui est évalué, et une perte de gamma le jour de la publication.
\textbf{19.} Un écart-type de 8.0~\% et un déplacement absolu attendu de 6.4~\%.
\textbf{20.} Parce qu'un événement ajoute de la variance à une date, et qu'une interpolation dans le temps la répartirait sur
des échéances qui se terminent avant lui.
\end{solution}

\begin{solution}{iq:dv:parametrising-the-surface:1}
$w(k)=a+b(\rho(k-m)+\sqrt{(k-m)^2+\varsigma^2})$~: $a$ relève tout le smile, $b$ ouvre l'angle entre
les ailes, $\rho$ la fait pivoter (skew), $m$ la décale à gauche ou à droite, $\varsigma$ arrondit le sommet.
\iqlookfor{la formule et le rôle géométrique de chaque paramètre.}
\end{solution}

\begin{solution}{iq:dv:parametrising-the-surface:2}
Calculer les variances totales~; la \omterm{def:dv:parametrising-the-surface:event}{variance d'événement} est la variance totale de l'échéance la plus tardive moins
le taux de variance de l'échéance antérieure multiplié par le temps de l'échéance la plus tardive~; sa racine carrée est l'écart-type du
mouvement, et son produit par $\sqrt{2/\pi}$ donne le déplacement absolu attendu.
\iqlookfor{raisonner en variance totale, et non en volatilité.}
\end{solution}

\begin{solution}{iq:dv:parametrising-the-surface:3}
La \omterm{thm:dv:parametrising-the-surface:lee}{formule des moments de Lee}~: la pente de la variance totale dans l'aile vaut $2-4(\sqrt{p^2+p}-p)$, où $p$ compte
les moments finis du sous-jacent~; elle vaut au plus 2, atteinte lorsqu'aucun moment au-delà du premier n'est
fini. Une aile plus pentue donnerait des prix de call qui violeraient les bornes d'arbitrage aux strikes extrêmes.
\iqlookfor{la borne et son lien avec les moments.}
\end{solution}

\begin{solution}{iq:dv:parametrising-the-surface:4}
Réduire la dimension~: à $(m,\varsigma)$ fixé, SVI est linéaire en trois paramètres, résolu par moindres carrés~;
chercher les deux paramètres restants à partir des valeurs de la veille (ou de la dernière seconde)~; imposer la borne de Lee et le
domaine admissible~; minimiser la distance hors des fourchettes avec des poids fonction de la largeur~; vérifier l'arbitrage après chaque
ajustement et revenir à la surface précédente si une vérification échoue.
\iqlookfor{réduction de dimension, démarrages à chaud, fourchettes, vérifications.}
\end{solution}

\begin{solution}{iq:dv:parametrising-the-surface:5}
Les valorisations du vendredi incluent deux jours de \omterm{def:dv:greeks-and-the-hedging-pnl:greeks}{theta} calendaire que le marché ne facture pas~; le lundi, les
options sont réévaluées à la hausse par rapport à ce que prévoit le modèle, et le desk enregistre un gain apparent de \omterm{def:dv:greeks-and-the-hedging-pnl:greeks}{vega} qui est
le \omterm{def:dv:greeks-and-the-hedging-pnl:greeks}{theta} du week-end qui revient. Il faut passer au \omterm{def:dv:parametrising-the-surface:business}{temps ouvré}, avec des pondérations pour le week-end.
\iqlookfor{l'horloge comme choix de modélisation.}
\end{solution}

\begin{solution}{iq:dv:parametrising-the-surface:6}
SSVI lorsque la surface doit être exempte d'arbitrage par construction et stable (systèmes de risque, valorisation d'options exotiques
échéances clairsemées ou illiquides), et pour extrapoler en échéance~; tranches indépendantes lorsque chaque échéance
doit coller de près à ses cotations (activité de teneur de marché). SSVI renonce à l'ajustement des smiles courts et pentus, car
un seul $\rho$ et une seule loi de courbure couvrent toutes les échéances.
\iqlookfor{garanties contre la flexibilité.}
\end{solution}
